% !TEX root = ../main.tex

\section{パンルヴェ方程式の基礎}

\subsection{変形微分方程式系}

%FIXME (証明とかを書き足す)
\subsection{2階フックス型微分方程式のラックス方程式}

\begin{equation}\label{eq1}
\frac{d^{2}z}{dx^{2}}=r(x,t)z
\end{equation}

\begin{equation}\label{eq2}
\frac{d\phi}{dx}+\frac{1}{2}p_{1}(x,t)\phi=0
\end{equation}

\begin{proposition}
    2階フックス型微分方程式\eqref{eq:2-ode}のモノドロミー保存変形と2階フックス型微分方程式\eqref{eq1}のモノドロミー保存変形とは、1階フックス型微分方程式\eqref{eq2}がモノドロミー保存変形を許すという条件のもとで同値である。
\end{proposition}

\subsection{変形微分方程式の解}

p.128
\begin{Q}[なぜ$n$階の変形微分方程式系を考えるのか？]
\textbf{A.} $\cdot$ $n=2$だけでやるよりも見通しがいいから。\\
\phantom{\textbf{A.} }$\cdot$ 線型微分方程式のモノドロミー保存変形によって非線形な関係が浮かび上がるから。
\end{Q}

フックス型微分方程式
\eqref{eq:linear-ode-n}を考える。\\
\begin{equation}\label{eq:auxiliary-phi}
\frac{d\phi(x)}{dx}+\frac{1}{n}p_{1}(x,t)\phi(x)=0
\end{equation}
で定まる解をとって、
\[
y=\phi(x,t)z
\]
とすると、\eqref{eq:fuchs_n_order}は
\begin{equation}\label{eq:reduced-fuchsian}
\frac{d^{n}z}{dx^{n}}+\sum_{j=2}^{n} q_{j}(x,t)\frac{d^{n-j}z}{dx^{n-j}} = 0
\end{equation}
となる。ここで、$q_{j}(x)$は$\phi(x,t)(i)\,(i=0,1,\cdots,j)$と$p_{i}(x,t)\,(i=1,\cdots,j)$の多項式である。\eqref{eq:reduced-fuchsian}のリーマンデータを
\[
(\mathbb{X}',\Xi',\rho')
\]
としておく。
ここで、\eqref{eq:reduced-fuchsian}の解の基本系を$\vec{\varphi}(x)$とする。解析接続をして、
\[
W(\vec{\varphi}^{\gamma})=W(\vec{\varphi})\Gamma(\gamma)
\]
命題\ref{prop:wronskian-ode}より、$w(\vec{\varphi})$は$x$によらないので、$w(\vec{\varphi}^{\gamma})=w(\vec{\varphi})$であるからdetをとって
\[
\det(\Gamma(\gamma))=1
\]
がわかる。よって、\eqref{eq:reduced-fuchsian}のモノドロミー群は$SL(n,\mathbb{C})$の部分群であることがわかった。すなわち、
\[
\rho':\pi(x_{0},\mathbb{X}')\to SL(n,\mathbb{C})
\]
である。ここで$\mathbb{X}'=\mathbb{P}^{1}(\mathbb{C})\setminus \Xi'$である。

\begin{definition}
    上のように、モノドロミー群が$SL(n,\mathbb{C})$の部分群であるような微分方程式を\textbf{SL型}という
\end{definition}

p.129
\begin{Q}[\eqref{eq:auxiliary-phi}に解はあるの？]
\textbf{A.} $\cdot$ $p_1(x,t)$が正則な$x$の周りではコーシーの存在定理（定理\ref{thm:cauchy}を参照）より必ず存在し、特異点の周りでも (複素べき)$\times$(正則関数とlogの多項式) の形で解を展開できる。(定理\ref{thm:fuchs}を参照)\\
\phantom{\textbf{A.} }$\cdot$ 【目的】この解 $\phi$ を使って元の変数を変換し、方程式から邪魔な $n-1$ 階微分の項を消し去る（正規化する）ため。
\end{Q}

\begin{Q}[リーマンデータって何？]
\textbf{A.} $\cdot$ 微分方程式の定義域、確定特異点の位置の集合、モノドロミーの組のこと。(定義\ref{def:riemann_data}を参照)\\
\end{Q}

\begin{Q}[解析接続と微分は交換可能？]
\textbf{A.} 可能。(定義\ref{prop:interchange-ac-diff}を参照)\\
\end{Q}

\begin{proposition}
    $n$階フックス型微分方程式\eqref{eq:linear-ode-n}のモノドロミー保存変形と、$n$階フックス型微分方程式\eqref{eq:reduced-fuchsian}のモノドロミー保存変形とは$1$階フックス型微分方程式\eqref{eq:auxiliary-phi}がモノドロミー保存変形を許すという条件のもとで同値である。
\end{proposition}%FIXME(証明したい)
%FIXME(モノドロミー保存変形とは)

%#endregion